\documentclass[11pt,a4paper]{article}
\usepackage{amsmath,amssymb,amsthm}
\usepackage{geometry}
\geometry{margin=1in}
\usepackage{hyperref}
\usepackage{enumitem}
\usepackage{xcolor}

\newtheorem{problem}{Problem}[section]
\newtheorem{definition}{Definition}
\newtheorem{theorem}{Theorem}
\newtheorem{proposition}{Proposition}

\title{A Rigorous Harmonic Sine Transform:\\
From Greedy Decomposition to Zeta‑Regularised Spectra\\
\large A Research Programme}
\author{}
\date{}

\begin{document}
\maketitle

\begin{abstract}
Victor Geere's greedy harmonic reconstruction of the sine function from the fractions \(\pi/(n+2)\) suggests a new non‑orthogonal basis for \(L^2[0,\pi]\). This programme develops a rigorous \emph{harmonic sine transform} (HST) based on that basis. The central objectives are: (1) to relate the HST to the classical Fourier transform and compare their properties; (2) to establish the HST as a well‑defined topological isomorphism onto a sequence space, with an explicitly computable dual basis; (3) to derive a Weil‑type explicit formula linking the zeros of the Riemann zeta function to HST coefficients; and (4) to construct and analytically continue zeta‑weighted Gram matrices whose spectral properties encode the non‑trivial zeros of \(\zeta(s)\). The outcome will be a new interdisciplinary theory at the interface of harmonic analysis, functional analysis, and analytic number theory, offering a novel operator‑theoretic perspective on the Hilbert–Pólya conjecture.
\end{abstract}

\section{Introduction and Motivation}
Victor Geere’s recent construction~\cite{Geere2026a,Geere2026b} shows that the sine function can be rebuilt from the \emph{harmonic fractions}
\[
\alpha_n = \frac{\pi}{n+2}, \qquad n=0,1,2,\dots,
\]
via a greedy step‑function algorithm. For any angle \(\theta\in[0,\pi]\) one defines recursively
\[
\theta_0 = 0,\qquad 
\delta_n(\theta) = \Theta(\theta-\theta_n-\alpha_n),\qquad
\theta_{n+1} = \theta_n + \delta_n(\theta)\,\alpha_n,
\]
where \(\Theta\) is the Heaviside step function. The sequence \(\sin(\theta_n)\) converges to \(\sin(\theta)\) with error \(O(1/n)\), and the selection indicators \(\delta_n(\theta)\) form a complete family of non‑orthogonal building blocks.

The present programme aims to turn this discovery into a fully fledged \emph{harmonic sine transform} (HST) – a linear operator on \(L^2[0,\pi]\) with a rich arithmetic structure – and to explore its deep connections with the Riemann zeta function \(\zeta(s)\). The four central objectives of the programme are:

\begin{enumerate}[label=(\arabic*)]
    \item Relate the HST to the classical Fourier transform and compare their analytic, numerical, and applicative properties.
    \item Extend the HST to a rigorous mathematical transformation, proving that it is a well‑defined topological isomorphism onto a sequence space, with a computable dual basis.
    \item Derive an explicit formula à la Riemann–Weil, linking the zeros of \(\zeta(s)\) to the HST of test functions via the von Mangoldt function.
    \item Establish analytic continuations for the associated zeta‑weighted Gram matrices and the Dirichlet series generated by HST coefficients, and identify their poles with the zeros of \(\zeta(s)\).
\end{enumerate}

Each objective is broken into mathematically well‑posed problems. Intermediate results are of independent interest and will be developed systematically.

\section{The Harmonic Sine Basis and Its Gram Matrix}

\subsection{Centred basis}
To obtain a zero‑mean family suitable for harmonic analysis we set
\[
e_n(\theta) = \delta_n(\theta) - \frac{\theta}{\pi},\qquad \theta\in[0,\pi].
\]
Each \(e_n\) is a piecewise‑constant function that changes its value at most once, on a set that depends on the whole greedy history up to index \(n\).  
The first task is to clarify the functional‑analytic status of \(\{e_n\}_{n=0}^\infty\).

\begin{problem}[Completeness and Riesz basis property]
\label{prob:riesz}
Prove that \(\operatorname{span}\{e_n\}\) is dense in \(L^2[0,\pi]\) and that \(\{e_n\}\) forms a \emph{Riesz basis} of this space. That is, there exist constants \(0 < A \le B < \infty\) such that for every finite sequence \(c_n\),
\[
A\sum_n |c_n|^2 \le \Big\|\sum_n c_n e_n\Big\|_{L^2}^2 \le B\sum_n |c_n|^2 .
\]
Equivalently, the Gram matrix
\[
K_{nm} = \langle e_n, e_m\rangle_{L^2[0,\pi]} = \int_0^\pi e_n(\theta)\,e_m(\theta)\,d\theta
\]
defines a bounded and boundedly invertible operator on \(\ell^2(\mathbb{N}_0)\).
\end{problem}

\begin{proof}[Strategy]
The greedy construction describes the support of \(e_n\) as a union of intervals with endpoints at certain rational multiples of \(\pi\). A detailed combinatorial analysis of the “harmonic digit expansion” is expected to yield an explicit asymptotic form for \(K_{nm}\), suggesting a near‑diagonal dominance that can be treated with classical perturbation arguments (e.g., Gershgorin circles for infinite matrices). Numerical experiments will precede a formal proof.
\end{proof}

\subsection{The dual basis}
A Riesz basis possesses a unique dual Riesz basis \(\{\tilde{e}_n\}\) with \(\langle e_n, \tilde{e}_m\rangle = \delta_{nm}\). Formally,
\[
\tilde{e}_n = \sum_m (K^{-1})_{nm} e_m .
\]

\begin{problem}[Explicit structure of the dual basis]
\label{prob:dual}
Compute the inverse Gram matrix \(K^{-1}\) explicitly, or at least its asymptotic decay. Show that the dual functions \(\tilde{e}_n\) are still piecewise‑constant with a controlled number of jumps, and that they can be computed efficiently by a recurrence that exploits the one‑step memory of the greedy process.
\end{problem}

A solution to Problem~\ref{prob:dual} is essential for a fast, numerically stable HST and for understanding the reconstruction formula.

\section{The Harmonic Sine Transform as a Rigorous Transformation}

\begin{definition}
For \(f\in L^2[0,\pi]\) we define its \emph{harmonic sine transform} (HST) by
\[
(\mathcal{H}f)_n = a_n = \langle f, e_n\rangle = \int_0^\pi f(\theta)\, e_n(\theta)\,d\theta .
\end{definition}

If Problem~\ref{prob:riesz} is resolved positively, \(\mathcal{H}\) maps \(L^2[0,\pi]\) linearly and continuously onto a weighted \(\ell^2\)‑space. The inverse is given by
\[
f = \sum_{n=0}^\infty a_n \tilde{e}_n,
\]
with convergence in \(L^2\) and, for sufficiently smooth \(f\), uniform. One may then speak of the \emph{harmonic sine series} of \(f\).

\paragraph{Comparison with Fourier sine series.}
The classical sine basis \(\{\sin(k\theta)\}_{k=1}^\infty\) is orthogonal with \(\|\sin(k\theta)\|^2 = \pi/2\). Their relationship with the \(e_n\) is encoded in the cross‑Gram matrix
\[
U_{k n} = \langle \sin(k\theta), e_n\rangle .
\]

\begin{problem}[HST–Fourier connection]
\label{prob:connect}
Compute \(U_{k n}\) exactly in terms of harmonic numbers. Use this to express the HST coefficients \(a_n\) as an infinite linear combination of Fourier sine coefficients \(b_k = \frac{2}{\pi}\int_0^\pi f(\theta) \sin(k\theta)\,d\theta\), and vice versa. Analyse the decay properties of the matrix \(U\) and interpret the HST as a \emph{non‑orthogonal time‑scale decomposition} whose “scale” index \(n\) corresponds to the harmonic fraction \(\pi/(n+2)\).
\end{problem}

The expected outcome is a precise dictionary between the two transforms, illustrating how the HST localises signal features on a harmonic scale that is intrinsically adapted to step‑wise approximations, whereas the Fourier transform uses globally supported sinusoids. This will open applications in signal processing where natural scales are harmonic (e.g., music, gear‑ratio analysis) and yield new approximation theorems.

\section{Relating the HST to the Riemann Zeta Function}

\subsection{HST Dirichlet series}

The harmonic fractions \(\alpha_n = \pi/(n+2)\) embed the denominators \(n+2\) directly into the basis. We define for any \(f\in L^2[0,\pi]\) the \emph{HST Dirichlet series}
\[
D_f(s) = \sum_{n=0}^\infty \frac{a_n}{(n+2)^s},\qquad a_n = \mathcal{H}f(n),
\]
initially for \(\operatorname{Re} s\) large enough. Because of the Riesz basis property, the sequence \((a_n)\) belongs to a weighted \(\ell^2\) space, implying absolute convergence of \(D_f(s)\) in some right half‑plane.

\begin{problem}[Meromorphic continuation of \(D_f\)]
\label{prob:contD}
For a dense class of test functions \(f\) (e.g., \(C^\infty[0,\pi]\) with vanishing moments at the endpoints), prove that \(D_f(s)\) extends meromorphically to \(\mathbb{C}\) and identify its poles. Conjecturally, the only poles are at \(s=1\) (coming from the constant term \(\int f(\theta)\frac{\theta}{\pi}d\theta\) multiplied by \(\zeta(s)-1\)) and possibly at the non‑trivial zeros of \(\zeta(s)\), if \(f\) is chosen appropriately.
\end{problem}

\begin{proof}[Strategy]
Express \(D_f(s)\) as an integral transform:
\[
D_f(s) = \int_0^\pi f(\theta) \Big( \sum_{n=0}^\infty \frac{e_n(\theta)}{(n+2)^s} \Big) d\theta
       = \int_0^\pi f(\theta) \Big( \sum_{n=0}^\infty \frac{\delta_n(\theta)}{(n+2)^s} \Big) d\theta 
         - \Big( \frac{1}{\pi}\int_0^\pi f(\theta)\theta\,d\theta \Big) (\zeta(s)-1).
\]
The inner sum \(\Phi(\theta,s) = \sum_{n=0}^\infty \frac{\delta_n(\theta)}{(n+2)^s}\) is a \emph{greedy harmonic zeta function} associated with the real number \(\theta/\pi\). Its analytic properties reflect those of \(\zeta(s)\) but are perturbed by the digit sequence \(\delta_n(\theta)\). Standard techniques from zeta functions of numeration systems (e.g., automatic sequences) will be adapted.
\end{proof}

\subsection{Zeta‑weighted Gram matrices}

A parallel construct is the family of operators \(K(s)\) on \(\ell^2\) defined by
\[
K_{nm}(s) = \int_0^\pi e_n(\theta) e_m(\theta) \Phi(\theta,s) \, d\theta .
\]

\begin{problem}[Spectral resolution of \(K(s)\)]
\label{prob:Kspec}
For \(\operatorname{Re} s>1\), \(K(s)\) is a trace‑class perturbation of a multiple of the Gram matrix \(K\). Show that the map \(s\mapsto K(s)\) extends meromorphically to the whole complex plane, and that its poles (or those of its Fredholm determinant) are precisely the zeros of \(\zeta(s)\). This would give a \textbf{harmonic‑sine analogue of the Hilbert–Pólya conjecture}: the zeros of \(\zeta(s)\) are the resonances of a naturally defined operator constructed solely from the harmonic sine decomposition.
\end{problem}

\section{An Explicit Formula Within the HST Framework}

The classical Weil explicit formula (for even test functions) reads
\[
\sum_{\rho} \widehat{h}(\rho) = h(0) + h(1) - \sum_{p}\frac{\log p}{p} \big( h(\log p) + h(-\log p) \big) + \ldots,
\]
where \(h\) is related to the test function via Mellin transform. We aim for a reformulation where the HST coefficients take the place of the Mellin transform.

\begin{problem}[HST Weil‑type identity]
\label{prob:weil}
Find a class of test functions \(f\in L^2[0,\pi]\) for which the sum over non‑trivial zeros \(\sum_{\rho} D_f(\rho)\) can be evaluated in closed form and expressed as a sum over prime powers involving the von Mangoldt function \(\Lambda(n)\). More precisely, establish an identity of the type
\[
\sum_{\rho} D_f(\rho) \;=\; \text{boundary terms} \;-\; \sum_{n=1}^\infty \Lambda(n)\, W_f(n),
\]
where \(W_f(n)\) is an explicit weight constructed from the HST of \(f\) (for instance, an average of \(f\) over certain intervals determined by the greedy expansion of \(\log n\)).
\end{problem}

\begin{proof}[How to proceed]
Start with the Perron formula for \(D_f(s)\):
\[
\frac{1}{2\pi i} \int_{c-i\infty}^{c+i\infty} D_f(s) \frac{x^s}{s} ds = \sum_{n+2 \le x} a_n + \text{error}.
\]
Choose \(x\) to be related to a prime power, and \(f\) so that its HST coefficients \(a_n\) mimic an arithmetic sequence. Because the Riesz basis is not orthogonal, this is an interpolation problem: for a desired sequence \((a_n)\in\ell^2\) we set \(f = \sum a_n \tilde{e}_n\). Thus, by choosing \((a_n)\) to be a smoothed version of the von Mangoldt function (or a related arithmetic sequence), we obtain a test function whose HST Dirichlet series is directly
\[
\sum_{n} \frac{\Lambda(n+2)}{(n+2)^s} = -\frac{\zeta'}{\zeta}(s) + \text{regular terms}.
\]
A whole family of explicit formulas follows by varying the sequence \((a_n)\) within the range of the HST.
\end{proof}

An intermediate, crucial step is a \emph{harmonic Poisson summation formula} for the HST: an identity equating a sum of HST coefficients to a sum of the function \(f\) evaluated at points related to the harmonic numbers. Such a formula would already mirror the symmetry underlying classical explicit formulas.

\section{Analytic Continuation of the HST‑Zeta Objects}

The final objective is to place the operator \(K(s)\) and the Dirichlet series \(D_f(s)\) on a firm analytic footing.

\begin{problem}[Global meromorphy of \(K(s)\)]
\label{prob:global}
Extend the holomorphic definition of \(K(s)\) from \(\operatorname{Re} s>1\) to a meromorphic family of compact operators on \(\ell^2\). Prove that \(\det(I - K(s))\) (a regularised Fredholm determinant) is an entire function whose zero set coincides with the multiset of non‑trivial zeros of \(\zeta(s)\), counting multiplicities. This would simultaneously give an operator‑theoretic characterisation of the Riemann zeta zeros and a new spectral interpretation of the Riemann Hypothesis as the reality of the eigenvalues of a self‑adjoint operator derived from \(K(1/2+it)\).
\end{problem}

\begin{proof}[Vector‑valued approach]
Treat the map \(s\mapsto \Phi(\cdot,s)\) as an \(L^2[0,\pi]\)-valued holomorphic function. Its Mellin‑like structure suggests that the continuation can be effected by expanding \(\Phi(\theta,s)\) into the dual basis \(\tilde{e}_n\) and using the known analytic continuation of the partial sums of the harmonic zeta function. The pole structure then emerges naturally from the singularity of \(\zeta(s)\) at \(s=1\) and from the logarithmic derivative \(-(\zeta'/\zeta)(s)\).
\end{proof}

\begin{problem}[Connection to the Hilbert–Pólya approach]
\label{prob:hp}
If \(K(1/2+it)\) is self‑adjoint for real \(t\), then its eigenvalues are real; by Problem~\ref{prob:global} these eigenvalues would encode the zeros of \(\zeta(1/2+it)\). This would realise the Hilbert–Pólya idea within the harmonic sine framework. Investigate the symmetry of the Gram matrix to find a canonical conjugation that makes \(K(1/2+it)\) self‑adjoint.
\end{problem}

\section{Comparison with the Fourier Transform and Applications}

Once the HST is rigorously established, a systematic comparison with the Fourier sine transform on \([0,\pi]\) will be carried out:

\begin{itemize}
    \item \textbf{Time‑scale localisation.} While the Fourier sine basis functions oscillate globally, each \(e_n\) is a step function that changes value near a specific cumulative harmonic sum. The HST therefore provides a “harmonic scalogram” – an angular analogue of a wavelet transform with scales \(\pi, \pi/2, \pi/3, \dots\).
    \item \textbf{Uncertainty principles.} Prove an analogue of the Heisenberg inequality for the HST, exploiting the non‑orthogonal Gram matrix to find pairs of sequences that minimise a joint concentration measure.
    \item \textbf{Numerical quadrature.} Exploit the fact that the HST of a function is computed by a finite number of threshold comparisons if the function is simple, leading to fast transforms for piecewise‑constant or piecewise‑polynomial data.
    \item \textbf{Arithmetic signal processing.} The harmonic fractions relate directly to the Riemann zeta function; the HST could serve as a natural tool for separating periodicities that are rationally related to the harmonic series – useful in analysing experimental data with logarithmic scaling.
\end{itemize}

\section{Conclusion and Logical Progression}

This programme transforms the speculative “harmonic sine pathway to the Riemann Hypothesis” into a sequence of concrete, mathematically well‑defined problems. The logical progression is:

\begin{enumerate}[start=1]
    \item \textbf{Foundations.} Solve the Riesz basis problem for \(\{e_n\}\); compute the dual basis and establish the HST as a continuous isomorphism. Publish the rigorous transform alongside the first numerical comparison with Fourier series.
    \item \textbf{Analytic continuation of \(D_f(s)\).} Analyse the primitive zeta function \(\sum \delta_n(\theta)/(n+2)^s\) and prove a meromorphic continuation for \(D_f(s)\). Derive an explicit formula for a restricted class of test functions whose HST coefficients are indicator sequences.
    \item \textbf{Operator‑theoretic zeta connection.} Extend the operator \(K(s)\) and prove that its zeros detect the Riemann zeros. Study the self‑adjointness of \(K(1/2+it)\) and its implications for the Hilbert–Pólya conjecture.
    \item \textbf{Synthesis.} Unify the greedy geometric construction with the spectral theory of \(\zeta(s)\), providing both new analytic tools and a conceptual bridge between harmonic analysis and number theory.
\end{enumerate}

Each stage is designed to produce significant results that are of independent value, even if the ultimate proof of the Riemann Hypothesis remains a long‑term goal. The approach offers a novel, mathematically sound route forward, grounded in the remarkable harmonic decomposition of the sine function.

\begin{thebibliography}{9}
\bibitem{Geere2026a}
V.~Geere, \emph{A Purely Geometric Reconstruction of the Sine Function},
March~2026. \url{https://victorgeere.co.za/maths/sine-harmonic.html}
\bibitem{Geere2026b}
V.~Geere, \emph{On the Correlation Structure of a Greedy Harmonic Decomposition of the Sine Function},
March~2026. \url{https://victorgeere.co.za/maths/greedy-harmonic-decomposition.html}
\end{thebibliography}

\end{document}